\documentclass[11pt,a4paper]{article}
\usepackage[margin=20mm,headheight=15pt]{geometry}
\usepackage[T1]{fontenc}
\usepackage{lmodern,microtype,amsmath,amssymb,booktabs,tabularx,longtable,graphicx}
\usepackage[table]{xcolor}
\usepackage{enumitem,fancyhdr,caption,placeins}
\usepackage[hidelinks]{hyperref}
\definecolor{ink}{HTML}{244F65}
\setlength{\parindent}{0pt}
\setlength{\parskip}{4pt}
\setlist{nosep,leftmargin=1.4em}
\captionsetup{font=small,labelfont=bf}
\pagestyle{fancy}
\fancyhf{}
\fancyhead[L]{\small CSE 516: Probabilistic Graphical Models}
\fancyhead[R]{\small M2 / Group 4}
\fancyfoot[C]{\small\thepage}
\newcommand{\E}{\mathbb E}
\newcommand{\Var}{\operatorname{Var}}
\newcommand{\N}{\mathcal N}
\newcommand{\figdir}{figures}
\newcommand{\code}[1]{\texttt{\small #1}}
\begin{document}
\begin{center}
{\small Ahmedabad University \quad | \quad Theoretical AI}\par
\vspace{5pt}
{\Large\bfseries Bayesian Personalized Recommendation using\\[3pt]
Latent Trait Models \& Expectation Propagation}\par
\vspace{7pt}
{\large Milestone 2: Technical Foundation}\par
\vspace{6pt}
\begin{tabular}{ll@{\qquad}ll}
AU2440170 & Jyotier Vithalani & AU2440234 & Tanmay Kanani\\
AU2440259 & Ansh Chaudhari & AU2440178 & Shlok Shah\\
AU2440221 & Het Patel & &
\end{tabular}\par
\vspace{6pt}
{\small Technical snapshot and experiment date: 17 October 2026}
\end{center}
\vspace{2pt}
\begin{abstract}
This milestone turns the M1 proposal into a reproducible study of approximate
Bayesian inference. We implement an identity-only, static Matchbox variant with
Gaussian latent traits, user/item biases and fixed ordinal rating thresholds.
The implementation exposes the distinction between single-pass assumed-density
filtering (ADF) and repeated expectation propagation (EP): EP stores and removes
old rating messages before replacing them. Product factors use variational
messages, so the algorithm is a hybrid EP/VMP procedure. On the full MovieLens
1M dataset, three paired random splits give RMSE $0.8980\pm0.0049$ for ADF and
$0.8633\pm0.0040$ after five EP passes at eight traits. Mean fitting time rises
from 2.64 to 13.63 seconds, and EP requires 219.77 MiB of rating-site storage.
The conventional matrix-factorization baseline reaches clipped-point RMSE
$0.8646\pm0.0031$, making its accuracy close to EP5 at lower measured cost.
Nine mathematical and implementation checks pass. A separate quadrature study
reveals variance and correlation errors in the approximation. These are initial
technical results, not a reproduction of the complete Matchbox system or a
claim of state-of-the-art performance.
\end{abstract}

\section{Problem, research question and progress from M1}
Sparse ratings leave substantial uncertainty about both user preferences and
item characteristics. A latent-trait model shares information between ratings
through common user and item variables. A Bayesian approximation also retains
parameter uncertainty, instead of storing only one fitted vector per entity.
Our theoretical question remains: \emph{what accuracy, probability quality,
computation and memory trade-offs arise when ADF is replaced by repeated EP?}

M0 selected Matchbox and MovieLens 1M. M1 specified the representation, baselines,
KPIs and experiments but contained no implementation or measured results. M2
now includes data validation, a concrete graph specification, parameter-learning
and prediction code, implemented baselines, numerical checks and recorded runs.
The graph structure is fixed; this project does not perform structure discovery.

\begin{table}[htbp]\centering\small
\begin{tabularx}{\linewidth}{@{}lXX@{}}\toprule
M1 direction & M2 implementation & Later work\\\midrule
ADF versus EP & Separate ADF path and cached-site EP; 1, 3, 5 passes & Longer runs, order and damping sensitivity\\
Latent dimensionality & $K=2,8$ on three splits & Wider dimension/prior study\\
Ordinal feedback & Fixed global cuts at 1.5, 2.5, 3.5, 4.5 & Learned user-specific thresholds\\
Metadata and sparsity & Metadata files validated; ID-only training & Metadata ablation; controlled user cold-start\\
SOTA positioning & Matchbox implementation anchor; PreRec conceptual comparison & Revisit after M3 experiments\\\bottomrule
\end{tabularx}
\caption{Scope changes are explicit. Deferred experiments are not reported as completed.}
\end{table}

\section{Probabilistic representation}
\subsection{Variables, observations and assumptions}
For users $u\in\{1,\ldots,U\}$ and items $i\in\{1,\ldots,I\}$, let
$s_u,t_i\in\mathbb R^K$ be latent traits and $b_u,b_i\in\mathbb R$ be biases.
For an observed pair $(u,i)\in\mathcal D$, define
\begin{align}
z_{uik}&=s_{uk}t_{ik}, &
a_{ui}&=\mu+b_u+b_i+\sum_{k=1}^{K}z_{uik},\label{eq:affinity}\\
y_{ui}\mid a_{ui}&\sim\N(a_{ui},\beta^2), &
R_{ui}=r&\iff c_{r-1}<y_{ui}\leq c_r.\label{eq:ordinal}
\end{align}
The rating is observed; traits, biases and the noisy affinity are latent.
The fixed threshold vector is $c=(-\infty,1.5,2.5,3.5,4.5,+\infty)$.
The plug-in offset $\mu$ is the training-set mean rating; its uncertainty is not
modelled. Priors factorize:
\begin{equation}
s_{uk}\sim\N(m^0_{uk},\sigma_s^2),\quad
t_{ik}\sim\N(n^0_{ik},\sigma_t^2),\quad
b_u,b_i\sim\N(0,\sigma_b^2).
\end{equation}
The small trait prior means are sampled once from $\N(0,0.1^2)$ using the run
seed and then held fixed. They break the bilinear sign symmetry without using
held-out labels. This is a modelling choice requiring sensitivity analysis,
not a parameter learned from the test data.

\subsection{Factorization and conditional independence}
Writing $\Theta=(S,T,b_U,b_I)$ and using the deterministic relations in
Equation~\ref{eq:affinity}, the joint distribution has the form
\begin{equation}
\begin{aligned}
p(\Theta,Z,A,Y,R)=p(\Theta)\prod_{(u,i)\in\mathcal D}
&\left[\prod_k\delta(z_{uik}-s_{uk}t_{ik})\right]
\delta\!\left(a_{ui}-\mu-b_u-b_i-\sum_k z_{uik}\right)\\
&\times\N(y_{ui};a_{ui},\beta^2)\,
\mathbf1\{c_{R_{ui}-1}<y_{ui}\le c_{R_{ui}}\}.
\end{aligned}\label{eq:joint}
\end{equation}
The Gaussian prior factors regularize traits and biases. Product factors express
preference compatibility; the sum factor combines compatibility and systematic
rating offsets. The Gaussian factor represents residual variability and the
interval factor maps a continuous score to an observed rating.

Ratings are conditionally independent given all relevant latent parameters and
fixed hyperparameters. Ratings sharing a user or item are generally dependent
after those parameters are marginalized out. Independent priors do \emph{not}
imply independent posterior parameters: the factorized Gaussian posterior is an
approximation. The model assumes stationary preferences and treats the observed
rating sample as its target population. It does not model why a user chose to
rate a movie or identify causal effects.

\begin{figure}[htbp]\centering
\includegraphics[width=.88\linewidth]{\figdir/factor_graph.pdf}
\caption{The implemented factorization for one rating. The dimension plate repeats
the product construction. User/item variables are shared by all corresponding
ratings. \code{graph.py} also exports an actual small bipartite graph to JSON;
the optimized numerical kernel implements these factors without allocating a
large Python object for every rating.}
\label{fig:graph}
\end{figure}

\section{Learning and inference}
\subsection{What is learned and what is queried?}
Training approximates $p(\Theta\mid R_{\mathcal D})$ by a product of univariate
Gaussians. The stored natural parameters are precision $\lambda=1/v$ and
precision-weighted mean $\eta=m/v$. Parameter learning therefore consists of
updating posterior beliefs about traits and biases. Thresholds, graph edges,
noise variance and prior variances remain fixed in the reference experiment.

Prediction asks for the full five-class distribution
$p(R_{ui}=r\mid R_{\mathcal D})$ on an unobserved rating. Let $m_s,v_s,m_t,v_t$
denote approximate trait moments. Independence within $q$ gives
\begin{align}
\E[s_kt_k]&=m_{s_k}m_{t_k},\\
\Var(s_kt_k)&=v_{s_k}v_{t_k}+v_{s_k}m_{t_k}^{2}+v_{t_k}m_{s_k}^{2},\label{eq:productvar}\\
m_y&=\mu+m_{b_u}+m_{b_i}+\sum_km_{s_k}m_{t_k},\\
v_y&=\beta^2+v_{b_u}+v_{b_i}+\sum_k\Var(s_kt_k).
\end{align}
These are exact moments under the factorized $q$; replacing the sum of products
by a Gaussian is an additional approximation. The resulting probability is
\begin{equation}
\widehat p_r=\Phi\left(\frac{c_r-m_y}{\sqrt{v_y}}\right)
-\Phi\left(\frac{c_{r-1}-m_y}{\sqrt{v_y}}\right).\label{eq:prediction}
\end{equation}
The main point prediction is $\widehat R=\sum_{r=1}^{5}r\widehat p_r$.
We also record the categorical median, the appropriate decision for absolute
error. The probability floor of $10^{-15}$ is applied only for numerical
protection and the probabilities are renormalized.

\subsection{Local hybrid EP/VMP update}
Matchbox uses variational messages for the bilinear product and EP or exact
Gaussian messages for other factors~\cite{stern2009}. Our implementation follows
that division rather than labelling every operation ``EP.'' During an update,
the user and item cavity distributions are kept fixed. A local loop alternates
forward product moments, Gaussian sum/noise propagation, ordinal moment matching
and reverse messages. The loop reconstructs each marginal from its cavity and
current message; it does not repeatedly add the same likelihood.

For the ordinal factor, let the incoming score distribution be $\N(m,v)$ and
set $\alpha=(c_{r-1}-m)/\sqrt v$, $\gamma=(c_r-m)/\sqrt v$ and
$Z=\Phi(\gamma)-\Phi(\alpha)$. The tilted moments are
\begin{align}
\widehat m&=m+\sqrt v\,\frac{\phi(\alpha)-\phi(\gamma)}{Z},\label{eq:truncmean}\\
\widehat v&=v\left[1+\frac{\alpha\phi(\alpha)-\gamma\phi(\gamma)}{Z}
-\left(\frac{\phi(\alpha)-\phi(\gamma)}{Z}\right)^2\right].\label{eq:truncvar}
\end{align}
Gaussian division produces its outgoing message:
$\lambda_{\mathrm{out}}=\widehat v^{-1}-v^{-1}$ and
$\eta_{\mathrm{out}}=\widehat m/\widehat v-m/v$.
Gaussian noise and sums then propagate the message toward each product.
If the reverse message to $z=st$ has moments $m_z^{\leftarrow},v_z^{\leftarrow}$,
the VMP message to $s$ is
\begin{equation}
\lambda_{f\to s}=\frac{\E[t^2]}{v_z^{\leftarrow}},\qquad
\eta_{f\to s}=\frac{m_z^{\leftarrow}\E[t]}{v_z^{\leftarrow}},\label{eq:vmp}
\end{equation}
with the analogous expression for $t$. The second moment
$\E[t^2]=v_t+m_t^2$ is essential. It is one reason VMP uncertainty can differ
from exact inference. Gaussian reverse messages update the biases.

The reference local loop uses at most 10 iterations, natural-parameter damping
0.7 and an absolute mean/variance change tolerance of $10^{-4}$. A capped local
update is retained and counted. Non-positive cavities, numerical underflow or
non-finite outputs are counted and skipped; no such skip occurred in the
reference benchmark. Appendix~\ref{app:derivation} details the reverse sums.

\subsection{ADF versus EP: avoiding double counting}
For each rating $n$, EP stores a factorized Gaussian site $\widetilde f_n$.
For a touched scalar parameter, the global approximation satisfies
\begin{equation}
\lambda=\lambda_0+\sum_n\lambda_n,\qquad
\eta=\eta_0+\sum_n\eta_n.\label{eq:site_sum}
\end{equation}
Before revisiting rating $n$, its old contribution is removed:
\begin{equation}
\lambda^{\setminus n}=\lambda-\lambda_n,\qquad
\eta^{\setminus n}=\eta-\eta_n.\label{eq:cavity}
\end{equation}
The local hybrid update returns a replacement site. On later passes its natural
parameters are damped with weight 0.5 toward the new site; the global parameters
become cavity plus replacement. The first pass uses weight 1.0. All sites begin
at zero, so the first EP pass equals ADF exactly under the same priors and order.

ADF immediately discards each site's messages after incorporating a new rating.
It needs no rating-site cache. Replaying the data while keeping the posterior
but omitting Equation~\ref{eq:cavity} would count observations again; it is not
the EP algorithm. The ADF interface explicitly rejects a second pass. The six
full-data ADF/EP1 comparisons have zero maximum natural-parameter difference.

There is no general convergence guarantee for this loopy, hybrid approximation.
Five outer passes are a declared computation budget, not a claim of convergence.
We log validation error, local iteration counts and parameter changes separately.

\section{Mathematical validation and computational cost}
\subsection{Independent checks}
The submitted validation command runs nine checks: ordinal moments against
independent numerical integration; an exact $K=0$ Gaussian/ordinal case;
ADF--EP1 equality; the prior-plus-sites identity; no double counting on revisiting
one rating; product moments against Monte Carlo; probability and metric sanity;
complete, disjoint, reproducible splits; and prediction-preserving checkpoint
reload. All nine passed. These checks validate specific calculations and
implementation invariants, not exact inference over the complete model.

\subsection{A bilinear case with a numerical reference}
To examine the approximation itself, use one rating $R=5$, offset 3 and noise
variance 1, with $s\sim\N(0.5,0.4)$ and $t\sim\N(0.8,0.4)$. Two zero-mean
biases have variance $10^{-8}$. Tensor-product Gauss--Hermite quadrature
integrates the likelihood $P(y>4.5\mid s,t)$ over the two traits, analytically
integrating the tiny biases. Increasing the order from 40 to 80 changes the
reported moments by at most $1.04\times10^{-10}$. This is a converged numerical
reference, not a symbolic proof. For this diagnostic only, the local loop is
allowed 500 iterations at tolerance $10^{-12}$ and converges after 50.

\begin{table}[htbp]\centering\small
\begin{tabular}{lrrrr}
\toprule
Quantity & Quadrature & Hybrid & Absolute error & Reference change \\
\midrule
$\mathbb E[s]$ & 0.922831 & 0.933888 & 0.011056 & 2.1e-11 \\
$\mathrm{Var}(s)$ & 0.325234 & 0.292029 & 0.033206 & 2.0e-11 \\
$\mathbb E[t]$ & 1.123960 & 1.136053 & 0.012093 & 3.9e-11 \\
$\mathrm{Var}(t)$ & 0.351319 & 0.315457 & 0.035862 & 1.0e-10 \\
$\mathrm{Cov}(s,t)$ & 0.070526 & 0.000000 & 0.070526 & 4.8e-11 \\
\bottomrule
\end{tabular}

\caption{One-observation validation. ``Reference change'' compares quadrature
orders 40 and 80. The nonzero exact posterior covariance cannot be retained by
the factorized approximation.}
\label{tab:quadrature}
\end{table}

The approximate means are within 0.013 of the reference, but both variances are
about 10\% too small. The covariance is forced to zero. Thus, passing local
formula checks is compatible with measurable model-level approximation error.
This is a concrete reason to investigate uncertainty quality in M3.

\subsection{Complexity and stored state}
For $N$ training ratings, $K$ traits, $L$ local iterations and $T$ outer passes,
the numerical work is $O(TNLK)$ (or $O(TNL(K+1))$ including the bias-only case).
Predicting $M$ ratings costs $O(MK)$. Posterior state uses
$O((U+I)(K+1))$ storage. Full EP additionally stores two natural parameters for
each of $2K+2$ variables touched by every rating:
\begin{equation}
\text{site bytes}=8\times2\times N(2K+2)=32N(K+1),
\end{equation}
using 64-bit floats. For $N=800{,}167$, $K=8$, this is 230,448,096 bytes, or
219.77 MiB. ADF's site storage is zero. At $K=8$, the trait/bias model has
$(6040+3706)\times9=87{,}714$ scalar latent parameters and stores 175,428
posterior natural-parameter values (about 1.34 MiB). Reported site memory is an
array-size calculation, not peak process RSS; data and temporary arrays add overhead.

\section{Dataset, baselines and experimental protocol}
\subsection{Preparation and split control}
MovieLens 1M is downloaded from GroupLens~\cite{movielens,harper2015}. The archive
and three data files are SHA-256 checked before use. The audit finds 1,000,209
ratings, 6,040 users, 3,706 distinct rated movies, and 3,883 movie metadata rows.
Movie catalogue size, maximum movie ID and number of rated movies are different
quantities. Density among the rated user/item vocabulary is 4.468\%, so the
matrix is approximately 95.532\% unobserved. No duplicate user--item pair is
present. Movie titles are read as Latin-1. Missing entries are unobserved,
not negative feedback or zero ratings.

Seeds 42, 43 and 44 independently permute all rating rows. Each split contains
800,167 training, 100,020 validation and 100,022 test ratings (flooring the
80\% and 10\% sizes leaves the remainder for test). Models use identical row
indices for a seed. ID vocabulary construction uses only IDs; means, histograms
and learned parameters use training ratings. Split indices and their hashes are
saved. ADF and EP process the same training order, and EP retains it on revisits.

No held-out rating has a user absent from training. The three test splits contain
26, 9 and 13 ratings, respectively, whose item is absent from training; these are
retained and use prior information. This small incidental set is not a designed
cold-start experiment. The split is random by interaction, not chronological,
and does not measure future deployment performance. The paper's user-level
cold-start protocol is different, so published Matchbox MAEs are not placed in
our numerical comparison table.

\subsection{Implemented baselines}
\textbf{Global, user and item means.} The global model uses training rating
frequencies $g_r$. For an entity with class counts $n_r$ and total $n$, the
smoothed categorical prediction is $(n_r+10g_r)/(n+10)$. Its expected rating is
the corresponding shrinkage mean; an unseen entity falls back to the global
distribution. This supplies proper probabilities without assigning a fake
posterior variance to a point estimator.

\textbf{Biased matrix factorization (MF).} We fit
$\mu+b_u+b_i+p_u^Tq_i$ by regularized squared-error SGD, with $K=8$, learning
rate 0.007, multiplicative decay 0.98 per epoch, regularization 0.03 and 25
epochs. Updates use the old user and item vectors for each rating. The epoch
with lowest validation clipped-point RMSE is retained (epoch 25 for all three
runs). This is an SVD-style latent-factor baseline, not a literal full-matrix
SVD or Bayesian inference method. A single Gaussian observation scale is fitted
on validation NLL within standard deviations $[0.2,3]$, then the same fixed
thresholds yield five probabilities. We report both the distribution-mean
decoder and the conventional clipped point prediction. The scale wrapper models
residual spread, not uncertainty over the MF parameters.

\textbf{ADF.} The single-pass path is the inference baseline for EP. Both paths
use exactly the same representation, priors and local update code. The two
trait counts, three reported EP checkpoints, and other settings were fixed before
the full reference test run. An earlier 100,000-training-row development run
checked execution; it is not part of the reported benchmark and did not change
the declared reference hyperparameters.

\begin{table}[htbp]\centering\small
\begin{tabularx}{\linewidth}{@{}lXlX@{}}\toprule
Setting & Value & Setting & Value\\\midrule
Trait counts & 2, 8 & Observation variance & $\beta^2=1$\\
Trait prior variance & 0.25 & Bias prior variance & 1\\
Prior mean scale & 0.1 & Local tolerance / cap & $10^{-4}$ / 10\\
Local damping & 0.7 & Outer damping & 1 on pass 1; 0.5 later\\
EP pass budget & 5 & Reported checkpoints & 1, 3, 5\\
Main prediction & Categorical expectation & Replicates & Seeds 42, 43, 44\\\bottomrule
\end{tabularx}
\caption{Reference configuration. No systematic hyperparameter optimization is claimed.}
\end{table}

\subsection{KPIs, timing and statistical treatment}
MAE and RMSE use the expected categorical rating unless explicitly labelled
``clipped point.'' NLL is the mean $-\log\widehat p_{R}$ in nats. The multiclass
Brier score is the mean sum of squared errors over the five class probabilities.
Top-class ECE uses 10 equal-width probability bins. Equal-tail 90\% rating
intervals and categorical-median MAE are also retained in the result CSV.

Tables report mean $\pm$ sample standard deviation over the three seeded runs.
The seeds change the split and initialization together; they do not isolate
these sources of variation. Their datasets overlap, so they are not three
independent samples from a population. We use paired descriptive differences
and make no significance claim from three runs. Full robustness and uncertainty
analysis belong to M3.

Runs used Python 3.12.14 on Linux x86-64, NumPy 2.3.5, SciPy 1.17.0 and Numba
0.68.0. The container exposed nine logical CPUs; the numerical kernels run
sequentially and thread environment variables were set to one. Kernel compilation
is warmed before timing. Fitting times exclude download, parsing and KPI
evaluation. MF time includes validation checkpoint selection and scale fitting.
Prediction time is recorded separately. These are wall times in a shared
container, not hardware-independent speed guarantees.

\section{Initial results and interpretation}
\subsection{Predictive comparison}
\begin{table}[htbp]\centering
\resizebox{\linewidth}{!}{\begin{tabular}{lrrrr}
\toprule
Method & MAE $\downarrow$ & RMSE $\downarrow$ & NLL $\downarrow$ & Fit time (s) \\
\midrule
Global mean & $0.9342\pm0.0018$ & $1.1174\pm0.0022$ & $1.4558\pm0.0018$ & $0.001\pm0.000$ \\
User mean & $0.8313\pm0.0019$ & $1.0353\pm0.0020$ & $1.3469\pm0.0014$ & $0.045\pm0.000$ \\
Item mean & $0.7827\pm0.0029$ & $0.9785\pm0.0037$ & $1.3205\pm0.0039$ & $0.045\pm0.005$ \\
MF, K=8 & $0.6907\pm0.0029$ & $0.8681\pm0.0031$ & $1.2210\pm0.0033$ & $3.000\pm0.565$ \\
ADF, K=8 & $0.7101\pm0.0039$ & $0.8980\pm0.0049$ & $1.2520\pm0.0046$ & $2.644\pm0.212$ \\
EP3, K=8 & $0.6887\pm0.0034$ & $0.8731\pm0.0046$ & $1.2271\pm0.0039$ & $8.187\pm0.421$ \\
EP5, K=8 & $0.6800\pm0.0029$ & $0.8633\pm0.0040$ & $1.2176\pm0.0032$ & $13.632\pm0.879$ \\
\midrule
MF, clipped point & $0.6809\pm0.0029$ & $0.8646\pm0.0031$ & same model & same model \\
\bottomrule
\end{tabular}
}
\caption{Test results on full MovieLens 1M, mean $\pm$ sample SD over three seeds.
MF's main row uses the probability-distribution mean, like the other rows.
The final row exposes its standard clipped point decoder; the same calibrated
probabilities supply its NLL. Times refer to fitting, not end-to-end runtime.}
\label{tab:main}
\end{table}

\begin{figure}[htbp]\centering
\includegraphics[width=\linewidth]{\figdir/baseline_comparison.pdf}
\caption{Prediction and probability quality. Error bars are sample SD across
three paired seeded runs, not confidence intervals. The plotted MF uses its
distribution-mean decoder.}
\end{figure}

At eight traits, EP5 reduces mean RMSE by 0.0347 relative to ADF (about 3.86\%)
and NLL by 0.0344 nats. Every seeded run shows the same direction. The most
plausible explanation is that later messages can incorporate information learned
from ratings processed after an earlier rating, while cavity removal avoids
counting that rating twice. This is an interpretation of the update mechanism
and observed trend, not a proof that more passes always help.

The comparison with MF is much closer. EP5's RMSE is 0.0048 below the MF
distribution-mean decoder, but only 0.0013 below the conventional clipped MF
predictor on average. The clipped MF predictor wins on seed 43. This difference
is too small and the run count too low to support a broad superiority claim.
MF is also faster in this implementation. Global/user/item means provide useful
checks that personalization improves over simple aggregate predictions.

\subsection{Iteration, dimension and resource trade-offs}
\begin{table}[htbp]\centering\small
\begin{tabular}{llrrrr}
\toprule
$K$ & Schedule & RMSE & NLL & ECE & Site memory (MiB) \\
\midrule
2 & ADF & 0.9011 & 1.2557 & 0.0408 & 0.00 \\
2 & EP3 & 0.8846 & 1.2386 & 0.0456 & 73.26 \\
2 & EP5 & 0.8760 & 1.2300 & 0.0482 & 73.26 \\
8 & ADF & 0.8980 & 1.2520 & 0.0449 & 0.00 \\
8 & EP3 & 0.8731 & 1.2271 & 0.0537 & 219.77 \\
8 & EP5 & 0.8633 & 1.2176 & 0.0562 & 219.77 \\
\bottomrule
\end{tabular}

\caption{Mean results by dimension and schedule. Complete per-run values and SDs
are included in \code{summary.csv}; site memory excludes parameter and data arrays.}
\label{tab:dimensions}
\end{table}

\begin{figure}[htbp]\centering
\includegraphics[width=\linewidth]{\figdir/ep_tradeoff.pdf}
\caption{Validation trajectory and test error versus cumulative fitting time.
Labels on the right are the prespecified EP pass counts. Error bars show SD.}
\label{fig:tradeoff}
\end{figure}

Eight traits improve the mean EP5 RMSE over two traits (0.8633 versus 0.8760),
but require three times as much site memory: 219.77 versus 73.26 MiB. These two
settings establish a measurable capacity trade-off; they do not locate an optimal
dimension. EP5 at $K=8$ costs about 5.16 times the measured ADF fit time.
Prediction remains about 0.56 microseconds per rating in the measured batch
because EP changes fitted parameters rather than the prediction graph.

No numerical update was skipped. The largest local-cap count is 1,027 out of
800,167 updates (0.128\%). By pass five at $K=8$, only 10--17 local updates per
seed reach the cap. Nevertheless, end-of-pass RMS changes in parameter means
remain 0.0328--0.0347 at that checkpoint. Five passes therefore do not establish
outer convergence. A local iteration count, an outer pass count and predictive
stabilization are different diagnostics.

\subsection{Probability quality does not reduce to RMSE}
At $K=8$, EP improves Brier score from 0.6681 to 0.6547, but top-class ECE rises
from 0.0449 to 0.0562. Lower NLL and better squared probability error do not
guarantee improvement in every calibration summary. The pooled reliability
curve indicates underconfident top-class predictions in substantial probability
regions. The 90\% equal-tail intervals cover about 98.22\% of test ratings and
span 3.82 rating categories on average for EP5. Discrete interval rounding makes
nominal coverage conservative, and fixed observation noise also affects this
spread. We do not call these intervals calibrated merely because coverage is high.

\begin{figure}[htbp]\centering
\includegraphics[width=\linewidth]{\figdir/calibration.pdf}
\caption{Reliability and calibration. Pooled bins are descriptive because the
three splits overlap. Bins with 50 or fewer pooled cases are omitted from the
curve only; all bins remain in the ECE computation and saved CSV.}
\label{fig:calibration}
\end{figure}

The local quadrature study's variance underestimation and the benchmark's broad
rating intervals are not contradictory. One concerns posterior uncertainty over
parameters in a small controlled problem; the other also includes fixed
observation noise, discretization and model mismatch. Global-mean ECE is very
small, yet its RMSE/NLL are poor: matching the frequency of one common class is
not evidence of useful personalized discrimination.

\section{Related work and revised SOTA positioning}
Stern, Herbrich and Graepel's Matchbox~\cite{stern2009}, developed at Microsoft
Research, is our methodological base. It combines latent feature mappings,
multiple feedback models and hybrid message passing; ADF and cached repeated
EP supply the central computational comparison. Our code deliberately restricts
that model to static ID features and fixed global thresholds. We neither use
Infer.NET nor reproduce its full message schedule, learned thresholds, metadata
or dynamics. Its published scores use different protocols and are not directly
comparable to ours.

Lin et al.'s PreRec~\cite{lin2024}, with AWS AI Labs affiliations, is the
WSDM 2024 anchor retained from M1. It uses a hierarchical Bayesian model with
neural user/item encoders, domain variables and popularity effects. Training
optimizes a MAP objective with Adam; zero-shot prediction alters the domain and
popularity mechanisms. Its XMarket/Online Retail experiments concern
cross-domain sequential recommendation and ranking metrics. Reported gains
support the value of modelling transfer and bias in that setting. They do not
demonstrate that EP is better or provide a MovieLens rating-prediction baseline.

\begin{table}[htbp]\centering\small
\begin{tabularx}{\linewidth}{@{}lXX@{}}\toprule
Aspect & This M2 study & PreRec anchor\\\midrule
Target & Predict 1--5 ratings within MovieLens & Transfer next-item preference across domains\\
Representation & Gaussian traits/biases; fixed ordinal cuts & Hierarchical domain, popularity and neural representations\\
Learning/inference & Factorized posterior; hybrid ADF/EP-VMP & MAP-trained latent/neural variables; intervention-based scores\\
Evidence & Paired rating splits; MAE, RMSE, NLL, time & Cross-market/platform ranking and limited-data adaptation\\
Open issue & Approximation and uncertainty quality & Dependence on representations and assumed bias mechanisms\\\bottomrule
\end{tabularx}
\caption{Conceptual comparison, not an empirical head-to-head result. PreRec is
the selected recent anchor from M1, not a claim to exhaust the 2026 literature.}
\end{table}

The positioning has therefore become narrower and more defensible since M1:
M2 establishes an inspectable inference benchmark and exposes its uncertainty
limitations. It does not establish causal debiasing, cross-domain transfer or a
new recommendation algorithm. PreRec motivates future questions about what
additional variables are needed; it is not used as an ornamental citation or
as an unrun numerical competitor.

\section{Limitations, M3 plan and conclusion}
The principal limits are fixed thresholds and noise, no metadata in training,
factorized Gaussian posteriors, approximate product messages, order-dependent
ADF, only two trait counts and three overlapping random splits. The MF and PGM
families have not received exhaustive or matched hyperparameter searches. Random
interaction evaluation can be optimistic for a future or new-user deployment.
No robustness, causal, cold-start or user-benefit claim follows from the current
results. Real timings also include implementation-specific advantages from
Numba-compiled loops.

For M3, we prioritize: (1) validation-only sensitivity studies for priors, noise,
damping and update order; (2) longer EP runs with stated stopping diagnostics;
(3) learned user thresholds and metadata with a controlled ablation; (4) true
user cold-start and reduced-training-density experiments; and (5) uncertainty
checks and resampling analysis that account for user clustering. These will
test which conclusions survive beyond this initial configuration.

M2 answers the required technical-foundation question: the representation,
learning, inference, baselines and KPI pipeline now run end to end on the stated
benchmark. EP improves on its ADF starting point at a measurable computation and
memory cost. Its accuracy is close to conventional MF, and its uncertainty
quality still needs work. The most useful outcome is a reproducible basis for
those M3 investigations, with the approximations and incomplete work visible.

\FloatBarrier
\bibliographystyle{plain}
\bibliography{references}

\appendix
\section{Reverse messages and executable algorithm}\label{app:derivation}
Suppose the ordinal factor supplies a Gaussian message
$\N(y;m_o,v_o)$. Reverse propagation through Gaussian noise gives variance
$v_o+\beta^2$. If the product output $z_k$ has forward moments $M_k,V_k$,
then the reverse sum message to that output is
\begin{align}
d_k &= m_o-\mu-m^{\setminus n}_{b_u}-m^{\setminus n}_{b_i}
                  -\sum_{\ell\ne k}M_\ell,\\
w_k &= v_o+\beta^2+v^{\setminus n}_{b_u}+v^{\setminus n}_{b_i}
                  +\sum_{\ell\ne k}V_\ell.
\end{align}
Substitute $m_z^{\leftarrow}=d_k$ and $v_z^{\leftarrow}=w_k$ in
Equation~\ref{eq:vmp}. For a bias, omit that bias's own cavity mean and variance
from the sum and retain all product moments. Its message precision is $1/w_b$
and weighted mean is $d_b/w_b$. Product forward messages use current local
marginals as required by the hybrid update; bias-to-sum inputs remain cavities.

The reason Gaussian division works is that
$\N(x;m,v)\propto\exp(-\lambda x^2/2+\eta x)$. Multiplication adds
$(\lambda,\eta)$; division subtracts them. The cache invariant in
Equation~\ref{eq:site_sum} is thus directly testable after any number of passes.

\begin{table}[htbp]\centering\small
\begin{tabularx}{\linewidth}{@{}rX@{}}\toprule
Step & One EP pass\\\midrule
1 & Visit training rows in the saved seeded order.\\
2 & Gather the user's and item's $K$ traits and two biases.\\
3 & Subtract this row's old natural-parameter sites from the global marginals.\\
4 & Run the local forward/backward hybrid update with the cavity fixed.\\
5 & Blend old and proposed sites using the specified outer damping.\\
6 & Set global marginals to cavity plus replacement, then overwrite the row's sites.\\
7 & Log numerical failures, local caps, parameter changes and validation KPIs.\\\bottomrule
\end{tabularx}
\caption{ADF uses the same local operation once per rating and discards the site cache.}
\end{table}

\section{Reproducibility and requirement coverage}
The GitHub-ready folder is named
\code{M2\_g4\_th\_apl\_recommend-sys\_team}. The root README and the
Code README give environment setup, quick-demo, validation, full-experiment,
figure-generation and report-build commands. Raw MovieLens data are not
redistributed: GroupLens' README requires separate permission for redistribution.
The downloader, source link and checksums are supplied instead.

\begin{longtable}{@{}p{.28\linewidth}p{.67\linewidth}@{}}
\toprule Requirement & Submitted evidence\\\midrule\endhead
Graph and distributions & Section 2; \code{graph.py}; \code{graph\_instance.json}; factor-graph figure\\
Parameter learning & \code{model.py}: \code{local\_messages}, \code{sweep}; no structure-learning claim\\
Inference output & \code{predict\_proba}; saved checkpoint; \code{demo.py}\\
Baseline and benchmark & \code{baselines.py}, \code{data.py}, \code{configs/m2.json}\\
Mathematical validation & Nine checks in \code{validation.py}; quadrature study and CSV\\
Initial KPIs & \code{metrics.csv}, \code{summary.csv}, plots and paired differences\\
Seeds and splits & Saved split indices, split hashes, cold-entity audit\\
Execution evidence & \code{run.log}, environment JSON, config, core-source hashes\\
Report source & This PDF, \code{main.tex}, BibTeX and generated tables\\
Video / presentation & To be supplied by the team, as requested; no completed video is claimed\\
Individual Google Form & Each member submits separately outside GitHub\\\bottomrule
\end{longtable}

The source audit used M1 commit
\code{16a72dc91e1acb3342c411ac264ceb9eeebe9856}. It contained the report,
draft diagram, READMEs and a Drive video link, but no presentation file or source
implementation. The linked video's content was not used as technical evidence.
The M2-specific notice specifies 15 October 2026, while the general guidelines
list 11 October; this report records the actual 17 October experiment date and
does not assert on-time submission or an extension.

\section{Team accountability and AI assistance}
The M1 responsibility plan is retained for review: Jyotier---mathematical model,
ADF/EP and repository coordination; Tanmay---baselines; Ansh---paper and SOTA
analysis; Shlok---dataset and metadata; Het---KPIs, experiment tracking and
documentation. These are inherited assignments, not verified records of who
completed the M2 work. Members must confirm their actual contributions in the
individual form; no contribution or signature has been invented here.

Generative AI assisted the implementation, debugging, experiment execution,
figures and report preparation. Numerical results were produced by the submitted
code, not invented from expectations or copied from either paper. The team is
responsible for reviewing the package, reproducing it and explaining its
mathematics, assumptions and limitations before submission. The accompanying
\code{TEAM\_REVIEW.md} and \code{DEFENSE\_NOTES.md} identify concrete review
tasks and explanations; they do not certify that this review has already happened.
\end{document}
